\chapter{Harmonic Oscillator Review, Spin-$\tfrac12$, and the First Exchange Principle}

These notes follow Leonard Susskind's fourth lecture in the supplementary \emph{Advanced Quantum Mechanics} sequence, with curation by LazyingArt LLC. The lecture unfolds in a single arc: a purposeful review of the harmonic oscillator as preparation for later field theory, a rapid but concrete return to spin-$\tfrac12$ as angular momentum, an atomic-physics motivation for spin and exclusion, and finally an intentionally unfinished opening of identical-particle statistics. Throughout the oscillator discussion Susskind sets \(m=1\) and mostly suppresses \(\hbar\); we follow that convention and only mention \(\hbar\) when the lecture does.

\section{Why the Harmonic Oscillator Returns}

Susskind begins by explaining why the harmonic oscillator keeps returning. Any system with an equilibrium configuration, once displaced slightly from equilibrium, will generally oscillate. A violin string is the easy example: it vibrates in a fundamental mode together with higher harmonics. But the real point is not acoustics. The oscillators that matter most for the rest of the course are the oscillatory modes of fields.

This is already visible in electromagnetism. The equilibrium configuration of the electromagnetic field is simply the vacuum: empty space with no electric or magnetic field. If one perturbs that equilibrium, for example by sending radiation into a cavity or waveguide, the field begins to oscillate in its allowed modes. Later those modes will be related to photons. The same pattern reappears in solids: a crystal does not consist of independent atoms oscillating separately, but of coupled atoms whose collective motions propagate as waves. Those oscillatory modes are again described by harmonic oscillators, and their quanta are phonons.

The review is therefore not a detour into familiar material. It is preparation for everything that follows: photons, phonons, and more generally the oscillatory structure that lies behind quantum field theory.

\section{Oscillator Algebra in Ladder Form}

With that motivation in place, Susskind returns to the algebraic description. In the convention \(m=1\), the Hamiltonian is
\begin{equation}
H=\frac{1}{2}p^{2}+\frac{1}{2}\omega^{2}x^{2},
\end{equation}
where \(\omega\) is both the oscillator frequency and the parameter setting the curvature of the potential.

The creation and annihilation operators are defined by
\begin{equation}
a_{\pm}=\frac{p\pm i\omega x}{\sqrt{2\omega}}.
\end{equation}
The factor \(\sqrt{2\omega}\) is not decorative. It is chosen so that the commutator takes the simplest possible form,
\begin{equation}
[a_{-},a_{+}]=1.
\end{equation}
The number operator is then
\begin{equation}
N=a_{+}a_{-}.
\end{equation}
This commutator, together with the definition of \(N\), implies that \(a_{+}\) and \(a_{-}\) raise and lower the eigenvalues of \(N\) by one unit. The spectrum of \(N\) is therefore
\[
0,1,2,3,\ldots
\]

The lecture also recalls how the Hamiltonian is expressed in terms of \(N\). Classically one would expect the product \(a_{+}a_{-}\) simply to reconstruct the Hamiltonian. Quantum mechanically the noncommutativity produces an extra term:
\begin{align}
N
&=
\frac{1}{2\omega}(p+i\omega x)(p-i\omega x) \\
&=
\frac{1}{2\omega}\bigl(p^{2}+\omega^{2}x^{2}+i\omega[x,p]\bigr) \\
&=
\frac{1}{2\omega}\bigl(p^{2}+\omega^{2}x^{2}-\omega\bigr),
\end{align}
since \([x,p]=i\) in the \(\hbar=1\) convention. Therefore
\begin{equation}
H=\omega\left(N+\frac{1}{2}\right).
\end{equation}
Susskind also writes the simplified form
\begin{equation}
H=\omega N
\end{equation}
when the additive constant is dropped. He is explicit about the distinction: the zero-point term \(\tfrac12\omega\) is real, but for many purposes it can be ignored because an additive constant in the energy does not change the equations of motion.

The extra \(\tfrac12\omega\) is the ground-state energy of the oscillator, or, in field-theory language, the vacuum energy. Susskind emphasizes that it is a consequence of the uncertainty principle: \(x\) and \(p\) cannot both vanish simultaneously, so the positive terms in the Hamiltonian cannot both be driven to zero.

\section{A Student Question, Wave Functions, and the Ground State}

At this point a student interrupts with a natural question: why does the ladder include \(0\) at all? Susskind first answers abstractly, before passing to wave functions. Because the Hamiltonian is built from \(p^{2}+\omega^{2}x^{2}\), it cannot have arbitrarily negative eigenvalues. There must therefore be a bottom state. The only way the ladder can stop is if the lowering operator annihilates that state:
\begin{equation}
a_{-}\lvert 0\rangle=0.
\end{equation}
Applying \(N=a_{+}a_{-}\) to the same state gives
\begin{equation}
N\lvert 0\rangle=0.
\end{equation}
So the bottom state has number eigenvalue \(0\). Susskind pauses over the notation: the symbol \(0\) here is merely the name of the state, not the zero vector.

Only after that abstract argument does he place the wave-function description ``parallel'' to the operator structure. State vectors become wave functions \(\psi(x)\), and energy eigenstates may be written as \(\psi_{0}(x),\psi_{1}(x),\ldots\). In the position representation,
\begin{equation}
X\mapsto x,\qquad P\mapsto -i\,\frac{d}{dx}.
\end{equation}
If one restores \(\hbar\), this becomes \(P\mapsto -i\hbar\,d/dx\); Susskind remarks that dimensional analysis tells one where the \(\hbar\)'s must go.

The time-independent Schr\"odinger equation is
\begin{equation}
\left[-\frac{1}{2}\frac{d^{2}}{dx^{2}}+\frac{1}{2}\omega^{2}x^{2}\right]\psi(x)=E\,\psi(x).
\end{equation}
As he stresses, this equation can be solved formally for any \(E\), but most formal solutions are not physical. For generic energy the wavefunction explodes at infinity and cannot be normalized. The admissible states are square-integrable wavefunctions,
\begin{equation}
\int_{-\infty}^{\infty}\psi^{*}(x)\psi(x)\,dx=1.
\end{equation}
Only for special values of \(E\) does the differential equation yield such honest eigenfunctions.

\paragraph{The smart route to \(\psi_{0}\).}
Rather than solve the second-order equation directly, Susskind uses the simpler first-order equation that must hold in the ground state:
\begin{equation}
a_{-}\psi_{0}(x)=0.
\end{equation}
Since an overall numerical factor is irrelevant in a homogeneous equation, one may write
\[
a_{-}\propto p-i\omega x.
\]
Substituting \(p=-i\,d/dx\) gives
\begin{equation}
(p-i\omega x)\psi_{0}(x)=0
\qquad\Longrightarrow\qquad
\left(\frac{d}{dx}+\omega x\right)\psi_{0}(x)=0.
\end{equation}
This is much simpler than the full Schr\"odinger equation. The standard trick is to set
\begin{equation}
\psi_{0}(x)=e^{f(x)}.
\end{equation}
Then
\[
\psi_{0}'(x)=f'(x)e^{f(x)},
\]
so the common factor \(e^{f(x)}\) drops out and one is left with
\begin{equation}
\frac{df}{dx}+\omega x=0.
\end{equation}
Integrating once,
\begin{equation}
f(x)=-\frac{1}{2}\omega x^{2}+C.
\end{equation}
Hence the ground-state wavefunction is
\begin{equation}
\psi_{0}(x)\propto e^{-\omega x^{2}/2}.
\end{equation}
This is the Gaussian Susskind draws on the board: smooth, concentrated near the origin, and rapidly decaying. It is manifestly square integrable. Its energy is already known from the ladder algebra,
\begin{equation}
E_{0}=\frac{\omega}{2}.
\end{equation}
If one wants the normalized form in the lecture convention, it is
\begin{equation}
\psi_{0}(x)=\left(\frac{\omega}{\pi}\right)^{1/4}e^{-\omega x^{2}/2}.
\end{equation}

\section{The First Excited State, Higher States, and the Semiclassical Picture}

Once the ground state is explicit, Susskind immediately asks for the first excited state. He again uses the algebra rather than attacking the differential equation afresh:
\begin{equation}
\psi_{1}(x)\propto a_{+}\psi_{0}(x).
\end{equation}
Since \(a_{+}\propto p+i\omega x\), one has
\begin{equation}
\psi_{1}(x)\propto (p+i\omega x)\psi_{0}(x).
\end{equation}
But from the ground-state equation \((p-i\omega x)\psi_{0}=0\), it follows that
\[
p\psi_{0}=i\omega x\,\psi_{0}.
\]
Therefore
\begin{equation}
(p+i\omega x)\psi_{0}=2i\omega x\,\psi_{0},
\end{equation}
and so, up to an irrelevant constant,
\begin{equation}
\psi_{1}(x)\propto x\,e^{-\omega x^{2}/2}.
\end{equation}

This is enough to read off the physics. The first excited state is an odd function:
\[
\psi_{1}(-x)=-\psi_{1}(x).
\]
Its probability density is even, but the wavefunction itself changes sign at the origin, and in fact vanishes there. Thus the probability of finding the oscillator exactly at the equilibrium point is zero in the first excited state. The derived form also makes clear that the first excited state has one node, at \(x=0\).

The lecture does not derive a closed formula for all higher states, and it is better not to force one in at this stage. The essential qualitative structure is already visible. Repeated action of \(a_{+}\) produces higher polynomials multiplying the same Gaussian envelope. Each higher state acquires another node; the parity alternates between even and odd; and the wavefunction is pushed farther outward from the origin. In standard notation these polynomials are Hermite polynomials, but the lecture's emphasis is not the name of the family. It is the picture.

That picture is simple and physical. More support at large \(|x|\) means more potential energy. More rapid oscillation means more momentum. As the energy increases, the wavefunction develops more and more ``wiggles'' and extends farther and farther into the wings. Susskind's point is that the oscillator is trying to do what classical intuition says it should do: high-energy states are likely to be found farther from the center, but when they pass through the center they carry larger momentum.

The lecture then pauses over the correspondence principle. A single high-\(n\) eigenstate is not itself a classical oscillator moving back and forth. To see something that looks classical, one must form a wave packet by superposing many energy eigenstates. Susskind's example is a displaced Gaussian: start with the ground-state packet and move it away from the origin. It is no longer an energy eigenstate, so under the time-dependent Schr\"odinger equation its components pick up different phases. The packet then swings back and forth, becoming more oscillatory near the center and smoother near the turning points, while maintaining its overall character. In that way the harmonic oscillator does satisfy the correspondence principle, but the classical motion is seen in a packet, not in a stationary eigenfunction.

The lecture ends the oscillator discussion with one further remark. Exactly solvable systems are rare. The harmonic oscillator is one of them, and the idealized hydrogen atom is another; both will continue to serve as laboratories for more general ideas.

\section{Spin-$\tfrac12$ as Angular Momentum}

Susskind then explicitly changes topic. The other subjects on his list for the evening are spin and the boson-fermion difference, and spin comes first. The return to spin is not meant as a fresh introduction. It is a reinterpretation of the earlier two-state system in the language of angular momentum.

For a particle at rest there is no orbital angular momentum \(\mathbf R\times \mathbf P\), but there may still be spin. One may picture that spin literally as a tiny object spinning, or one may regard it as an abstract intrinsic property attached to the particle. Susskind is deliberately permissive about the picture. What matters is the mathematics: whatever this extra degree of freedom is, it transforms as angular momentum.

The decisive criterion is the commutation relations. If three operators are attached to the spatial directions \(x\), \(y\), and \(z\), and if they satisfy the angular-momentum algebra, then they are angular momentum. In particular,
\begin{equation}
[L_{z},L_{x}]=iL_{y},
\end{equation}
with the other relations obtained by cyclic permutation.

The candidate operators are the Pauli matrices,
\begin{equation}
\sigma_{z}=
\begin{pmatrix}
1&0\\
0&-1
\end{pmatrix},
\qquad
\sigma_{x}=
\begin{pmatrix}
0&1\\
1&0
\end{pmatrix},
\qquad
\sigma_{y}=
\begin{pmatrix}
0&-i\\
i&0
\end{pmatrix}.
\end{equation}
Susskind chooses the representation in which \(\sigma_{z}\) is diagonal. Its eigenvalues are \(+1\) and \(-1\), corresponding to spin up and spin down along the \(z\)-axis.

The board calculation proceeds by explicit multiplication:
\begin{align}
\sigma_{z}\sigma_{x}
&=
\begin{pmatrix}
1&0\\
0&-1
\end{pmatrix}
\begin{pmatrix}
0&1\\
1&0
\end{pmatrix}
=
\begin{pmatrix}
0&1\\
-1&0
\end{pmatrix}
=
i\sigma_{y},\\
\sigma_{x}\sigma_{z}
&=
\begin{pmatrix}
0&1\\
1&0
\end{pmatrix}
\begin{pmatrix}
1&0\\
0&-1
\end{pmatrix}
=
\begin{pmatrix}
0&-1\\
1&0
\end{pmatrix}
=
-i\sigma_{y}.
\end{align}
Therefore
\begin{equation}
[\sigma_{z},\sigma_{x}]=2i\sigma_{y}.
\end{equation}
This is almost the angular-momentum algebra, but not quite. The mismatch is the factor of \(2\). The fix is exactly the one visible on the board:
\begin{equation}
S_{i}=\frac{\sigma_{i}}{2}.
\end{equation}
Then the operators become
\begin{equation}
\frac{\sigma_{z}}{2}=
\begin{pmatrix}
\tfrac12&0\\
0&-\tfrac12
\end{pmatrix},
\qquad
\frac{\sigma_{x}}{2}=
\begin{pmatrix}
0&\tfrac12\\
\tfrac12&0
\end{pmatrix},
\qquad
\frac{\sigma_{y}}{2}=
\begin{pmatrix}
0&-\tfrac{i}{2}\\
\tfrac{i}{2}&0
\end{pmatrix},
\end{equation}
and now
\begin{equation}
[S_{z},S_{x}]=iS_{y},
\end{equation}
again with cyclic permutations. The two-state system is therefore not merely analogous to angular momentum; it is the simplest nontrivial angular-momentum system.

Its \(z\)-component has eigenvalues
\begin{equation}
m_{s}=\pm\frac12.
\end{equation}
This is the spin-$\tfrac12$ doublet. Susskind emphasizes that many of the particles one most naturally thinks of as fundamental --- electrons, muons, neutrinos, quarks --- carry precisely this kind of spin.

\section{From Spin to Atoms and Shell Counting}

Once the algebraic point is in place, the lecture turns to why nature should care. There are two kinds of angular momentum: orbital angular momentum and spin. Their vector sum is the total angular momentum,
\begin{equation}
\mathbf J=\mathbf L+\mathbf S.
\end{equation}
For a particle,
\[
\mathbf L=\mathbf R\times \mathbf P.
\]
As Susskind later stresses, angular momentum is always defined relative to a chosen origin. In atomic physics the natural origin is the nucleus. The detailed rules for adding \(\mathbf L\) and \(\mathbf S\) are postponed; the lecture only needs the notation and the physical distinction.

The atomic motivation begins with the hydrogen atom. In the lecture convention,
\begin{equation}
L^{2}=\ell(\ell+1),
\end{equation}
and for fixed \(\ell\) the multiplicity is
\begin{equation}
2\ell+1.
\end{equation}
Susskind redraws the spectrum schematically as an energy-versus-\(\ell\) plot. The actual hydrogen energies are negative, but for the drawing he simply shifts the picture upward and focuses on degeneracies rather than scale.

\begin{figure}
\centering
\begin{tikzpicture}[x=1.25cm,y=0.95cm]
\draw[->] (-0.75,0) -- (-0.75,5.35);
\draw[->] (-0.95,0) -- (5.8,0);
\node at (-1.0,5.08) {$E$};
\node at (5.55,-0.28) {$\ell$};

\node at (0,-0.35) {$0$};
\node at (2.4,-0.35) {$1$};
\node at (4.8,-0.35) {$2$};

\foreach \y in {0.55,1.45,2.20,2.90,3.50,4.00,4.40,4.75}
  \draw (-0.18,\y) -- (0.18,\y);

\foreach \y in {1.45,2.20,2.90,3.50}
{
  \draw (2.02,\y) -- (2.17,\y);
  \draw (2.27,\y) -- (2.42,\y);
  \draw (2.52,\y) -- (2.67,\y);
}

\foreach \y in {2.20,2.90,3.50}
{
  \draw (4.22,\y) -- (4.33,\y);
  \draw (4.43,\y) -- (4.54,\y);
  \draw (4.64,\y) -- (4.75,\y);
  \draw (4.85,\y) -- (4.96,\y);
  \draw (5.06,\y) -- (5.17,\y);
}

\node[anchor=west] at (0.55,4.95) {$L^{2}=\ell(\ell+1)$};
\node[anchor=west] at (0.55,4.48) {$2\ell+1$};
\end{tikzpicture}
\caption{Schematic hydrogenic degeneracy pattern, redrawn from the lecture board and transcript. The vertical spacing is not to scale; the important content is the alignment of the \(\ell=1\) and \(\ell=2\) ladders with higher \(\ell=0\) states.}
\end{figure}

The \(\ell=0\) column contains the \(s\)-wave tower. The first \(\ell=1\) level is degenerate with the second \(\ell=0\) level, the first \(\ell=2\) level with the third \(\ell=0\) level, and so on. Counting the states shell by shell gives
\begin{equation}
1,\qquad 1+3=4,\qquad 1+3+5=9,\qquad 1+3+5+7=16,
\end{equation}
or more compactly,
\begin{equation}
1,4,9,16,\ldots=n^{2}.
\end{equation}
This is only the orbital counting. Once spin is included, the idealized shell totals double to
\begin{equation}
2,8,18,32,\ldots = 2n^{2}.
\end{equation}

The lecture then backs up and turns that counting into a physical argument. Helium is well described, at first approximation, by putting two electrons into the lowest orbital state. But when one proceeds to lithium, the third electron does not drop into the same state. It enters the first excited shell. Why should that happen, when it would seem energetically cheaper to place all the electrons in the lowest orbital?

Pauli's answer is the exclusion principle: no two electrons can occupy the same state. But that raises an immediate problem, because helium plainly allows two electrons in the same orbital. The way out is to recognize that the electron has more structure than its orbital motion alone. It carries an additional two-valued property. In the lecture this is introduced in the simplest possible language: call the two values up and down. Then two electrons may occupy the same orbital state provided they have opposite spin. They are not, after all, in the same full quantum state.

This is the lecture's route to the second row of the periodic table. The first excited shell contains one \(\ell=0\) state and three \(\ell=1\) states, for a total of four spatial states. Doubling by the two spin values gives eight possibilities, matching the eight entries in the second row of the periodic table in this very idealized shell picture.

Susskind also mentions two further points that keep the discussion honest. First, spectroscopy pointed in the same direction as chemistry, although the periodic-table argument is more transparent. Second, once the new two-valued property is interpreted as angular momentum, it ought to reveal itself in a magnetic field. A charged particle carrying angular momentum behaves like a tiny magnet, so the two spin orientations should split in energy when a magnetic field is applied. That splitting is indeed observed.

At the same time, the lecture is explicit that this shell model is naive. It ignores electron-electron interactions and therefore breaks down rather quickly for larger atoms. The simple \(2n^{2}\) pattern is useful here as motivation for spin and exclusion, not as a quantitatively complete many-electron theory.

\section{Identical Particles and the First Exchange Principle}

This brings Susskind to the second advertised subject: bosons and fermions. A student asks whether the exclusion principle is a separate postulate. Susskind's answer is historically layered. In nonrelativistic quantum mechanics it can be stated as a postulate. In relativistic quantum theory it becomes deeper. But for the moment he wants a cleaner quantum-mechanical formulation, one phrased directly in terms of identical particles.

He begins by stepping back to classical statistical mechanics. Suppose two particles of the same kind are placed in boxes. If one imagines that the particles secretly carry names --- Harry and Sally is the lecture's example --- then placing Harry in box 1 and Sally in box 2 seems different from placing Sally in box 1 and Harry in box 2. One may nevertheless choose to count those as the same configuration, and in classical statistical mechanics that choice leads to the familiar \(1/n!\) factors. The point is that classically one can imagine faint labels, even if no experiment can detect them.

Quantum mechanics is stricter. One does not have the luxury of treating truly identical particles as secretly labeled. Swapping two electrons is an honest transformation on the quantum state, not a bookkeeping convention.

For one particle the wavefunction is \(\psi(x)\). For two identical particles the state is described by a wavefunction of two arguments,
\begin{equation}
\psi(x_{1},x_{2}),
\end{equation}
where \(x_{1}\) and \(x_{2}\) label the two particle slots, not two spatial directions. The probability density \(|\psi(x_{1},x_{2})|^{2}\) is the probability of finding one particle at \(x_{1}\) and the other at \(x_{2}\).

The next move is the key one. Consider the operation that swaps the two particles. Susskind denotes it by \(S\); to avoid confusion with the spin operator \(\mathbf S\), we shall write it as \(P_{12}\). Its action is
\begin{equation}
P_{12}\psi(x_{1},x_{2})=\psi(x_{2},x_{1}).
\end{equation}
A function of two variables need not be symmetric, so in general this produces a new wavefunction.

Now swap twice. Nothing mysterious happens: interchanging the two particle slots a second time returns the original state. Therefore
\begin{equation}
P_{12}^{2}=1.
\end{equation}
Susskind also notes that exchange is a transformation on the Hilbert space and so must be unitary. A unitary operator whose square is the identity can only have eigenvalues
\begin{equation}
\lambda=\pm 1.
\end{equation}
Thus the physically relevant two-particle states fall into two classes:
\begin{equation}
P_{12}\psi=\psi,
\qquad \text{or} \qquad
P_{12}\psi=-\psi.
\end{equation}

That is exactly where the lecture stops. The wavefunction of two identical particles either returns to itself under exchange or returns with a minus sign. The full interpretation of those two possibilities as bosonic and fermionic statistics is deferred to the next lecture. Susskind is about to begin with the \(+1\) case when he announces that he is fading and leaves the rest for next time.

\section{Summary}

The lecture begins with the ubiquity of the harmonic oscillator and with a clear warning that this review is preparation for later field theory, not repetition for its own sake. From there it rebuilds the oscillator algebra, explains the zero-point shift, and answers a student question about why the ladder bottoms out. The abstract argument for a lowest state is then matched by an explicit position-space construction of the Gaussian ground state and the odd first excited state.

Only after that oscillator machinery is in place does the lecture pivot to spin, showing that the Pauli matrices become angular momentum operators once normalized as \(S_{i}=\sigma_{i}/2\). The atomic discussion then uses hydrogenic degeneracies and naive shell counting to motivate Pauli's extra two-valued electron property, later recognized as spin. Finally, exclusion is reformulated as the first statement of identical-particle exchange: for two identical particles, swapping the particle slots leaves the wavefunction either unchanged or multiplied by \(-1\). That deliberately unfinished dichotomy is the chapter's endpoint and the next lecture's starting point.